\begin{center}
    \large

    \begin{singlespace}
        \textbf{\englishTitle{}} \\[0.5cm]
    \end{singlespace}

    \begin{singlespace}
        Student : \studentEnName{}  \hspace{1.0cm}
        Advisor: Dr.\, \advisorEnName \\
    \end{singlespace}

    \vspace{0.5cm}
    \begin{singlespace}
        \NameofDepartmentInstituteEN\\
        National Yang Ming Chiao Tung University\\[0.2cm]
    \end{singlespace}

    \textbf{Abstract} \\[0.5cm]

\end{center}

\normalsize

Large vision-language-action (VLA) robot policies can map language and visual observations to continuous actions, but their inference cost makes per-control-tick replanning expensive. Action chunking amortizes the slow planner call by predicting multiple future actions at once, yet the later actions in a chunk may be executed under stale observations. This thesis studies whether a small wrist-camera residual module can improve frozen \smolvla{} action chunks without fine-tuning the slow planner.

\method{} uses \texttt{HuggingFaceVLA/smolvla\_libero} as a frozen slow planner and records an HFRVLA LeRobotDataset v3 containing base actions, slow-planner context, chunk indices, robot states, expert actions, and frozen DINOv3 wrist patch features. The trainable fast path predicts a 7D residual, and deployment uses a bounded merge, $\afinal=\abase+\alpha\cdot\clip(\deltah,-\dmax,\dmax)$. This design keeps global task interpretation in the frozen planner while injecting current local visual evidence from the wrist camera.

The main experiments focus on LIBERO-Spatial and are tracked by an evaluation registry. In a synchronous plan=50 execution/replan sweep, \method{} achieves $79\%,80\%,71\%,75\%,59\%,53\%$ success for $K=2,4,8,16,32,50$, compared with \smolvla{} at $72\%,65\%,60\%,57\%,50\%,40\%$, while underperforming at $K=1$ ($75\%$ vs. $79\%$). A synchronous matched chunk sweep shows gains for most $K\geq4$ settings, but long chunks remain difficult. Residual calibration confirms that residual authority is a first-order deployment variable: in the plan=50, exec=8 calibration, $\alpha=0.75,\delta_{\max}=0.2$ reaches $41/50=82\%$. The async-timestep planner-delay stress test is reported separately and shows that the fast residual path can reduce degradation relative to a disable-fast wrapper when planner chunks arrive late.

The thesis concludes that \method{} is best understood as a bounded local correction layer for frozen VLA action chunks, not as a replacement for frequent replanning, full asynchronous scheduling, or real-robot validation. Multi-seed evaluation, causal visual ablations, same-backbone A2C2 comparison, and physical robot validation remain future work.

\vspace{1cm}

\textbf{Keywords: vision-language-action model, action chunking, robotic manipulation, wrist camera, residual correction, DINO features, LeRobot.}
